\documentclass[10pt,a4paper,twocolumn]{article}

% ===== XeLaTeX + Korean =====
\usepackage{kotex}
\usepackage{fontspec}
\setmainfont{Times New Roman}
\setmainhangulfont{AppleMyungjo}
\setsanshangulfont{Apple SD Gothic Neo}

% ===== Layout =====
\usepackage[a4paper,top=18mm,bottom=20mm,left=14mm,right=14mm,columnsep=8mm]{geometry}
\usepackage{fancyhdr}
\setlength{\headheight}{24pt}
\usepackage{enumitem}
\usepackage{graphicx}
\usepackage{ragged2e}

\setlength{\columnseprule}{0.3pt}
\setlength{\parindent}{1em}
\linespread{1.12}

% ===== Header / Footer =====
\pagestyle{fancy}
\fancyhf{}
\renewcommand{\headrulewidth}{0pt}
\fancyhead[C]{\fontsize{20}{20}\selectfont\sffamily 영어 영역}
\fancyhead[LE,RO]{\fontsize{20}{20}\selectfont\bfseries\thepage}
\fancyhead[RE,LO]{\framebox[12mm][c]{\sffamily 고1}}
\renewcommand{\footrulewidth}{0pt}
\fancyfoot[C]{\framebox[18mm][c]{\thepage\,/\,1}}

% ===== Problem helpers =====
\newcommand{\probno}[1]{\noindent{\fontsize{12}{14}\selectfont\bfseries #1.}\ }
\newcommand{\instr}[1]{{\sffamily #1}\par\vspace{2pt}}

\newlist{choices}{enumerate}{1}
\setlist[choices]{label=,leftmargin=1.4em,itemsep=1pt,topsep=2pt,parsep=0pt}

\usepackage{mdframed}
\newmdenv[linewidth=0.5pt,innertopmargin=5pt,innerbottommargin=5pt,
          innerleftmargin=6pt,innerrightmargin=6pt,skipabove=4pt,skipbelow=4pt]{pbox}

\begin{document}

% ===================== 18 =====================
\probno{18}\instr{다음 글의 목적으로 가장 적절한 것은?}
\begin{pbox}
Dear Members,\par
\hspace{1em}I am Howard Bennett, the head of the Greenfield Hiking Club. As you know, our club has enjoyed the trails of Mount Carlton for many years. Recently, however, we have noticed that several sections of the path have been damaged by heavy rain, leaving them unsafe for hikers. To keep our trails enjoyable for everyone, we are organizing a trail repair day next month. We are looking for members who can spare a few hours to clear fallen branches and rebuild worn steps. No special skills are required—only a willingness to help. If you can join us, please sign up on our website by the end of this week. Your participation will make a real difference.\par
Sincerely,\par
Howard Bennett
\end{pbox}
\begin{choices}
\item ① 등산 안전 수칙을 안내하려고
\item ② 등산 동호회 가입을 권유하려고
\item ③ 등산로 폐쇄 일정을 공지하려고
\item ④ 등산로 보수 작업 참여를 요청하려고
\item ⑤ 등산 장비 기부를 부탁하려고
\end{choices}
\vspace{6pt}

% ===================== 19 =====================
\probno{19}\instr{다음 글에 드러난 Sora의 심경 변화로 가장 적절한 것은?}
\par\noindent Sora had spent weeks rehearsing for the school orchestra audition. As she waited backstage, her hands trembled and her heart pounded against her chest. She could hardly breathe when her name was finally called. Walking onto the stage, she worried that all her practice might go to waste. But the moment she placed her bow on the strings, the familiar melody flowed out smoothly, just as she had rehearsed it countless times. When she finished, the judges nodded and smiled warmly at her. Sora walked off the stage with a wide smile, glad that her long hours of effort had finally paid off.
\begin{choices}
\item ① bored $\rightarrow$ excited
\item ② nervous $\rightarrow$ proud
\item ③ confident $\rightarrow$ embarrassed
\item ④ hopeful $\rightarrow$ disappointed
\item ⑤ curious $\rightarrow$ frightened
\end{choices}
\vspace{6pt}

% ===================== 20 =====================
\probno{20}\instr{다음 글에서 필자가 주장하는 바로 가장 적절한 것은?}
\par\noindent Many parents rush to fix every difficulty their children face. When a child forgets her homework, the parent drives it to school; when two children argue, the parent steps in at once. These actions come from love, but they can rob children of valuable lessons. Small failures and conflicts are exactly how children learn to think, plan, and recover on their own. A child who is allowed to face the consequences of a forgotten assignment learns responsibility far better than one who is always rescued. Of course, parents should shield children from real danger. But for everyday troubles, stepping back is often the wiser choice. Give your children room to stumble, and they will grow stronger for it.
\begin{choices}
\item ① 자녀의 안전을 최우선으로 고려해야 한다.
\item ② 자녀에게 명확한 행동 규칙을 정해 주어야 한다.
\item ③ 자녀와 함께 보내는 시간을 늘려야 한다.
\item ④ 자녀의 실수를 꾸짖기보다 칭찬해야 한다.
\item ⑤ 사소한 문제는 자녀가 스스로 해결하도록 두어야 한다.
\end{choices}
\vspace{6pt}

% ===================== 21 =====================
\probno{21}\instr{밑줄 친 \textit{carry our own weather}가 다음 글에서 의미하는 바로 가장 적절한 것은? [3점]}
\par\noindent We often blame our surroundings for our moods. A rainy morning, a crowded train, a careless remark—these become the reasons we feel the way we do. Yet the same gray sky can leave one person gloomy and another perfectly at peace. The difference lies not outside but within. Those who have grasped this truth \underline{carry their own weather}. They walk into a tense room and lower the temperature with their calm; they meet bad news and choose their response rather than surrender to it. This does not mean ignoring reality or forcing a smile. It means recognizing that, while we cannot control every storm around us, we hold the thermostat of our inner climate. The skies outside may darken, but the forecast inside remains ours to set.
\begin{choices}
\item ① predict difficulties before they arrive
\item ② reshape the environment to suit our needs
\item ③ determine our inner state regardless of circumstances
\item ④ adapt quickly to sudden changes in life
\item ⑤ share our true emotions openly with others
\end{choices}
\vspace{6pt}

% ===================== 22 =====================
\probno{22}\instr{다음 글의 요지로 가장 적절한 것은?}
\par\noindent We tend to treat every empty moment as wasted time. Waiting for a bus, we reach for our phones; resting between tasks, we feel guilty and hunt for something productive to do. But the mind needs idle time as much as the body needs sleep. When we stop actively working on a problem and simply let our thoughts wander, the brain quietly keeps making connections in the background. Many great ideas arrive not at the desk but in the shower, on a walk, or while gazing out a window. By filling every gap with stimulation, we leave no room for these quiet insights to surface. Sometimes the most productive thing we can do is to allow ourselves to do nothing at all.
\begin{choices}
\item ① 규칙적인 운동은 두뇌 활동을 활발하게 한다.
\item ② 충분한 여유 시간은 창의적인 생각이 떠오르는 바탕이 된다.
\item ③ 목표를 세우면 시간을 효율적으로 쓸 수 있다.
\item ④ 집중을 위해서는 주변의 방해 요소를 없애야 한다.
\item ⑤ 디지털 기기 사용을 줄이면 수면의 질이 높아진다.
\end{choices}

\end{document}
